\section{Purchasing the nonlinear estimates on one scale sequence}
\label{sec:nonlinear-payment-selection}
The nonlinear estimates above are uniform in the finite torus once
the block size, derivative orders and stage parameters have been
fixed. We now choose coupling thresholds that pay them together.
In particular, the marked and unmarked constructions can use the
same transition. The two logarithmic corrections absorbed into the
forced gas require two successive payments; both are included.

\subsection{An elementary exponential payment}
Let $\beta_f(x)=x+c_L$, with fixed $c_L>0$, let
$B(x)=\lceil(\log x)^2\rceil$, and let $P_f(x)$ be the shifted
fine monitor. For $C>0$ put
\[
             \tau_C(x)=\exp\{-CP_f(x)^2/2\}.
\]
Fix $D_0>0$ and a function $f$ of negative upper exponent. For all
sufficiently large $x$, the following assertion holds simultaneously
for every pair $a,b\ge0$ such that
\[
             a\le D_0 e^{-CP_f(x)^2},\qquad b\le f(x).
\]
There is $\kappa\ge0$ with
\begin{equation}\label{eq:nonlinear-payment}
 \begin{gathered}
 b\le1,\qquad (2B+1)^4 I\le1,\qquad
 e^{-\kappa}(e^I-1)\le2\tau_C(x),\\
                    I=e^\kappa a+2b.
 \end{gathered}
\end{equation}
Take $\kappa=CP_f(x)^2/2$. Then
$e^\kappa a\le D_0\tau_C(x)$, so
\[
          (2B+1)^4I\le(2B+1)^4(D_0\tau_C+2f).
\]
The right side tends to zero. Indeed the geometric factor has
upper exponent zero, $f$ has negative upper exponent, and
$\tau_C$ has every negative upper exponent by the monitor estimate.
The same choice of threshold gives $b\le1$ and the middle inequality
of \eqref{eq:nonlinear-payment}. Since $(2B+1)^4\ge1$, also
$0\le I\le1$. Convexity of the exponential on $[0,1]$ gives
$e^I-1\le(e-1)I\le2I\le2$. Multiplication by
$e^{-\kappa}=\tau_C$ proves the last inequality.
These are exactly the smallness and output-price conditions for
the forced-gas absorption in Section~\ref{sec:soft-log-fusion}.
The threshold depends on $C,D_0,f,L$, but not on $a,b$.

\subsection{The two soft correction prices}
We verify the negative power needed in the second payment.
Fix a stage, including its derivative order, spatial weight,
support exponent, path widths and path prices. Let $t\in(0,1]$
be its threshold and $B$ its radius. Let $\alpha$ be the coupling
reset, with $|\alpha|\le A$, where $A\ge0$ is fixed after the
block length has been chosen. If the regular remainder allowance
is $0\le H\le t^2$, the reset's boundary-hit coefficient is
\[
       h_{\rm hit}=2t\left(|\alpha|+H/t^2\right)
                                      +\tfrac12 H C_{\rm pair}.
\]
Here $C_{\rm pair}\ge0$ depends only on the fixed derivative
orders. Since $H\le t^2\le t$,
\[
       h_{\rm hit}\le[2(A+1)+C_{\rm pair}/2]t.
\]

There are two correction families, with the support reserve used
for the first absorption and the larger reserve used for the second.
For either family let $R_0,K_0\ge0$ be respectively its reset and
path prices computed at radius one. The rooted footprint sum gives
\[
 R(B,H)\le R_0(1+B)^4H,\qquad
                         K(B)\le K_0(1+B)^4.
\]
To see the radius dependence explicitly, its footprint population
is a fixed nonnegative multiple of either $(2B+1)^4$ or
$(8B+5)^4$. The inequalities
$2B+1\le3(1+B)$ and $8B+5\le13(1+B)$ compare them with their
radius-one values. All remaining factors are fixed: exponential
support charges, inverse positive width powers and the convergent
geometric series for the unused width reserve. The same calculation
works with the larger reserve, using its smaller positive width.
Thus the total soft price satisfies
\begin{equation}\label{eq:nonlinear-soft-price}
 R(B,H)+h_{\rm hit}K(B)
 \le C_A(1+B)^4t,\qquad
 C_A=R_0+[2(A+1)+C_{\rm pair}/2]K_0.
\end{equation}
These constants are selected separately for the two families.

At coupling $x$, take
\[
 t(x)=x^{-1/2}\exp((\log\log x)^2),\qquad
 H(x)=c_{\rm st}x^{-11/10},\qquad 0<c_{\rm st}\le1/64.
\]
For all sufficiently large $x$, we have $t(x)\le1$.
For $x>1$ we also have $H\le x^{-1}\le t^2$, so the preceding
calculation applies eventually. Moreover, for any fixed constant $C_A$,
\[
 C_A(1+B(x))^4t(x)\le\beta_f(x)^{-1/4}
\]
eventually. Divide the left side by $x^{-1/4}$ and take logarithms:
the variable terms are
$4\log(1+\lceil(\log x)^2\rceil)+(\log\log x)^2-
\tfrac14\log x$, which tend to minus infinity. Finally
$\beta_f(x)/x\to1$, so the same assertion holds with the shifted
scale. This proves the required soft bound uniformly for
$|\alpha|\le A$, both derivative budgets and both finite-volume
families.

The first hard price of the nonlinear expansion is at most
$\exp(-\mu P_f^2/4)$, with $\mu>0$. Its accompanying soft price
is at most $\beta_f^{-1/8}$. In the multiple-insertion version it
is at most $C_{\rm cov}\beta_f^{-1/8}$, where the positive constant
$C_{\rm cov}$ depends only on the fixed enlargement depth. Both
have negative upper exponent. Apply
\eqref{eq:nonlinear-payment} with $C=\mu/4$, $D_0=1$ to obtain
the first output bound
\[
                     2e^{-\mu P_f^2/8}.
\]
Use this as the hard price for the second correction, whose soft
price is bounded by \eqref{eq:nonlinear-soft-price}. A second
application with $C=\mu/8$, $D_0=2$ gives
\begin{equation}\label{eq:nonlinear-two-payments}
          2e^{-\mu P_f^2/8}\quad\longmapsto\quad
                         2e^{-\mu P_f^2/16}.
\end{equation}
These are the two successive forced allowances used in the
nonlinear transition, including the larger-reserve correction.
All inequalities use the same value of $x$.

\subsection{Choosing thresholds without changing the transition}
The marked construction first chooses a sufficiently large block
size to pay its footprint conversion. For that fixed block and
fixed stage orders, Sections~\ref{sec:marked-activity-jets}
and~\ref{sec:leaf-expansion} give a threshold beyond which the
local marked derivatives, the two connected expansions and their
tail estimates hold. The chart transport radius is chosen before
this threshold; the microscopic fluctuation threshold tends to
zero and eventually fits inside it. The Gaussian and all-species
estimates are uniform in the finite torus and in the insertion,
so these choices do not depend on either of them.

There are finitely many permitted target stages at each current
stage $j$. For each such pair $(j,i)$, choose the block-size floor
given by the marked construction and take their finite maximum.
Choose the actual block size $L_j$ above that floor. Only after
$L_j$ is fixed do we choose the coupling thresholds needed for
\eqref{eq:nonlinear-payment}, \eqref{eq:nonlinear-soft-price}
and the fixed-order marked estimates. This order permits their
constants to depend on the actual block length and reset bound.

More explicitly, let $d_j$ be a previously chosen entry threshold,
and let $r_{j,i}$ be the threshold for all the requirements at
transition $(j,i)$. Replace it by
\begin{equation}\label{eq:nonlinear-entry-max}
 d'_j=\max\left(d_j,\max_{i\ {\rm permitted\ at}\ j}r_{j,i}\right).
\end{equation}
For any later thresholds $d''_j\ge d'_j$ and every $x\ge d''_j$,
all the requirements still hold. This is immediate from
$x\ge d'_j\ge r_{j,i}$; the other selected data, including the
block sizes and the transition maps, are unchanged. Applying
\eqref{eq:nonlinear-entry-max} successively to the ordinary and
the larger-reserve multiple-insertion requirements gives one set
of thresholds satisfying both. No uniform maximum over all
orders or all stages is needed.

Consequently the numerical arrival conditions previously required
by constructions with one or several insertions and by the
centered construction can be imposed simultaneously on the same selected
scale sequence. This is a purchase of their actual displayed
estimates. The construction of the underlying funded Gaussian
data and the subsequent coefficient and normalization identities
remain the respective analytic and algebraic steps already
separated in those arguments.
